%% ============================================================
%%  Front matter — Motivation, Objectives and Structure
%%
%%  Author      : Carlos David Prado-Socorro
%%  Institution : ICMol, Universitat de València
%%
%%  Thesis-level roadmap placed before Chapter 1. It states the
%%  problem, the device idea, the evidence hierarchy, and the
%%  six-chapter structure. Compiles standalone or \include'd
%%  from thesis.tex (guarded by \thesismode, as the chapters are).
%% ============================================================

%% ---- Standalone preamble (skipped when \include'd from thesis.tex) ----
\ifdefined\thesismode\else
  \documentclass[12pt, twoside]{report}
  \usepackage{thesis-format}
  \addbibresource{../bibliography/references.bib}
  \begin{document}
\fi

%% ============================================================
\chapter*{Motivation, Objectives and Structure of the Thesis}
\addcontentsline{toc}{chapter}{Motivation, Objectives and Structure of the Thesis}
\label{sec:scope_objectives}
%% ============================================================

\markboth{Motivation, Objectives and Structure}{Motivation, Objectives and Structure}

The cost of computing is increasingly the cost of \emph{moving data}, not the cost of arithmetic. In a conventional computer, the processor and the memory are physically separate, and every operation requires operands to be fetched across the channel that joins them. On modern hardware, a single off-chip memory access can cost three to four orders of magnitude more energy than the arithmetic it feeds~\cite{Horowitz2014}, and the workloads that now dominate---large machine-learning models and continuous streams from physiological and environmental sensors---are precisely those that demand enormous numbers of such accesses~\cite{LeCun2015}. As classical transistor scaling no longer hides this imbalance, the energy and latency of data movement have become the limiting factor of digital computing~\cite{Hennessy2019}. This thesis is motivated by that limit, and by one of the physical strategies proposed to relax it: building computation out of devices whose own internal dynamics carry and transform information, so that signalling, memory, and processing share a single physical substrate, as they do in the biological synapse~\cite{Mead1990,Tanaka2019}.

The device idea pursued here is the \emph{volatile organic memristive device}: a simple two-terminal element, made from an ion-containing polymer composite, whose conductance retains a decaying record of its recent electrical history. Rather than treating this volatility as a defect to be engineered away (as it would be for a non-volatile memory cell), the thesis treats the fading-memory timescale as the functional resource. A device that forgets at a controlled rate is a physical \emph{leaky integrator}, and a collection of such devices with different timescales is a substrate for \emph{temporal} (reservoir-style) computing on time-varying signals. The central claim of the work is that the chemistry of the polymer composite---above all its composition---provides a handle on that timescale, and that the resulting dynamics can be put to computational use.

The thesis builds this argument on a graded evidence hierarchy, stated plainly at the outset:

\begin{itemize}
  \item \textbf{A fully characterised exemplar (Chapter~2).} A single proof-of-concept Super~Yellow/Hybrane/lithium-triflate device is characterised with the complete synaptic measurement suite---current--voltage hysteresis, potentiation and depression, excitatory post-synaptic current, separated short- and long-term retention, and spike-timing-dependent plasticity---and given a mechanistic interpretation in terms of cation--anion mobility asymmetry.
  \item \textbf{A campaign-scale reproducibility and provenance layer (Chapter~3).} The subsequent Hybrane reproducibility collapse is diagnosed across the fabrication archive, separating finished-device shelf stability from batch-scale supply reproducibility and motivating the switch to PEO. The same provenance record and feature pipeline then become the methodological basis for the quantitative chapters.
  \item \textbf{A replicated composition spine with a sample-limited chemical landscape (Chapter~4).} A poly(ethylene oxide)-based corpus is built under a common protocol. Its quantitative spine is a \emph{replicated} composition series; host, anion, and cation substitutions are reported as an explicitly sample-limited chemical-tuning landscape, not as powered material laws. In particular, the \ce{Li+}\,>\,\ce{Na+}\,>\,\ce{K+} cation ordering is treated as a hypothesis to be tested, not assumed.
  \item \textbf{An in-silico temporal-computing demonstration from measured dynamics (Chapter~5).} Compact behavioural models extracted from the Chapter~4 measurements are used in reservoir-computing simulations, with heterogeneity claimed only where the simulations support it and negative controls (including the affective-label null on real physiological streams) reported explicitly.
\end{itemize}

\section*{Open Questions and Objectives}
\label{subsec:open_questions}

The work is organised around three groups of questions, each answered by a stage of the evidence hierarchy.

\emph{First, realising and understanding an exemplar.} Can a fully reversible two-terminal organic memristive device be made from a solution-processed polymer-electrolyte composite using standard organic-electronic fabrication? Can one device support analogue switching, distinguishable short- and long-term memory, an excitatory post-synaptic current response, and spike-timing-dependent plasticity? Can the fast and slow memory components be explained by the composite's cation--anion mobility asymmetry?

\emph{Second, composition as the quantitative lever, with chemistry as a sample-limited landscape.} Within a common host and a common alkali-metal triflate family, how do the device dynamics depend on composition when the polymer-to-semiconductor and salt-to-semiconductor mass fractions are varied around a reference formulation? This composition dependence is the axis the thesis sets out to \emph{replicate}. Host, anion, and cation substitutions are posed only as sample-limited hypotheses about the same dynamics---in particular, whether moving the cation along the \ce{Li+}, \ce{Na+}, \ce{K+} series shifts the measured fading-memory timescale as coordination chemistry suggests, a question the work tests rather than presumes.

\emph{Third, utility for temporal computing.} Can the volatile response of these devices be exploited for temporal computation? Can composition supply task-matched fading-memory timescales, while sample-limited chemistry and device-to-device spread populate a heterogeneous bank? And can measured variability contribute to computation rather than merely degrade it?

These questions translate into six ordered objectives:

\begin{enumerate}
  \item \label{subsec:objectives} Design, fabricate, and characterise a proof-of-concept two-terminal organic memristive device based on a Super~Yellow, Hybrane, and lithium-triflate composite, processed from a single solution onto a standard ITO/composite/Ag stack, and demonstrate the basic memristive and synaptic functionality in a single platform.
  \item Provide a mechanistic account of the device's two memory regimes in terms of the displacement of mobile triflate anions and lithium cations within the oxygen-rich ion-conducting host, supported by hard--soft acid--base reasoning on the cation--oxygen interaction.
  \item Build a comparative experimental corpus on a poly(ethylene oxide)-based platform loaded with alkali-metal triflate salts, fabricated under a common protocol, and characterise it through three common dynamical measurements: I--V hysteresis, variable-number-of-pulses potentiation, and variable-delay depotentiation.
  \item Quantify the composition dependence through a replicated PEO and lithium-triflate concentration series, and report host, anion, and cation substitutions as an explicitly sample-limited chemical-tuning landscape on the same three measurements.
  \item Construct experimentally grounded behavioural models from those measurements, using composition-indexed fading-memory times, nonlinear pulse-update rules, device-to-device variability envelopes, and read-transfer behaviour.
  \item Demonstrate, in simulation, temporal-computing functionality from these models in a focused reservoir-computing framework: memory-capacity and NARMA benchmarks, WESAD physiological-context reconstruction, and WESAD affective-computing tasks, with heterogeneity claimed where supported and negative controls reported explicitly.
\end{enumerate}

\section*{Structure of the Thesis}
\label{subsec:outline}

The thesis is organised in six chapters that follow this objective list.

\paragraph{Chapter~1 --- Introduction.} Provides the conceptual and materials background, moving from the limits of conventional computing, through the biological synapse and the memristor, to the polymer-electrolyte composite adopted as the device strategy for the rest of the thesis.

\paragraph{Chapter~2 --- Proof of Concept.} Reports the Super~Yellow, Hybrane, and lithium-triflate device. It is the only device in the thesis for which the complete synaptic suite is treated as primary evidence, and it isolates the cation--host coordination chemistry that the comparative work then tests. The chapter is built on Prado-Socorro et al.~\cite{PradoSocorro2022}.

\paragraph{Chapter~3 --- Reproducibility, Degradation, and the Device-Provenance Methodology.} Documents the reproducibility crisis that ended the Hybrane fabrication line and motivated the switch to a PEO host. Under explicit controls (per-device weighting, a standard-protocol corpus, sweep-amplitude matching), it establishes the batch-over-time collapse of the Hybrane switching window, offers a moisture-driven hydrolytic-aging hypothesis for the host stock, and introduces the fabrication-provenance record, the normalized characterization procedure, and the automated feature pipeline that the remaining chapters rely on.

\paragraph{Chapter~4 --- Compositional and Chemical Control.} Reformulates the proof-of-concept architecture on a PEO host loaded with alkali-metal triflates. Its quantitative spine is the replicated PEO/lithium-triflate composition grid; host, anion, and cation substitutions are reported as an illustrative, sample-limited chemical-tuning landscape rather than a powered Li/Na/K law. All claims rely on the same three common measurements.

\paragraph{Chapter~5 --- Data-Driven Temporal Computing.} Turns the experimental measurements into temporal-computing functionality. Compact behavioural models extracted from the Chapter~4 dynamics drive reservoir-computing simulations: memory-capacity and NARMA-style benchmarks, WESAD physiological-context reconstruction, and WESAD affective-computing demonstrations, with heterogeneity claimed where supported and nulls reported explicitly.

\paragraph{Chapter~6 --- Conclusions and Outlook.} Brings together the device-chemistry axis and the application axis, benchmarks the work against temporal, reservoir, and event-driven neuromorphic hardware, and states the limitations of the evidence base.

%% ---- Standalone close (skipped when \include'd from thesis.tex) ----
\ifdefined\thesismode
  \let\thesisstandaloneclose\relax
\else
  \printbibliography[heading=bibintoc, title={References}]
  \def\thesisstandaloneclose{\end{document}}
\fi
\thesisstandaloneclose
